\section*{अध्याय \ref{ch:g10:layered-ocean} --- परतदार महासागर}

\begin{solution}{exo:g10:layered-ocean:1}
\omterm{def:g10:layered-ocean:mixed-layer}{मिश्रित परत} महासागर की वह सतही परत है जिसमें तापमान और
लवणता गहराई के साथ लगभग नहीं बदलते; \omterm{def:g10:layered-ocean:thermocline}{ताप-प्रवणता परत} उसके नीचे की वह परत है जिसमें
तापमान गहराई के साथ तेज़ी से गिरता है। पवन और लहरें
\omterm{def:g10:layered-ocean:mixed-layer}{मिश्रित परत} को मथती हैं (सर्दियों में सतह का ठंडा होना भी इसमें सहायता करता है)।
\end{solution}

\begin{solution}{exo:g10:layered-ocean:2}
\omterm{def:g10:layered-ocean:upwelling}{उत्स्रवण} सतही परत के नीचे से ठंडे जल का ऊपर उठना है, जो पवन द्वारा
दूर धकेले गए सतही जल का स्थान लेता है। ऊपर उठता जल खनिज
पोषक तत्त्वों को सूर्य-प्रकाशित परत में लाता है, जहाँ पादप-प्लवक बढ़ने के लिए उनका उपयोग करते हैं;
पादप-प्लवक बाक़ी खाद्य शृंखला को भोजन देते हैं।
\end{solution}

\begin{solution}{exo:g10:layered-ocean:3}
उष्णकटिबंधीय प्रशांत: सतह पर लगभग \qty{29}{\degreeCelsius} और
\qty{1000}{m} पर \qty{4.5}{\degreeCelsius}। आर्कटिक महासागर: सतह पर लगभग
\qty{-1.4}{\degreeCelsius} और \qty{1000}{m} पर
\qty{0}{\degreeCelsius} के निकट।
\end{solution}

\begin{solution}{exo:g10:layered-ocean:4}
\qty{0}{\degreeCelsius} पर (\qty{1028.1}{kg/m^3}, जबकि
\qty{25}{\degreeCelsius} पर \qty{1023.3}{kg/m^3})। समुद्री जल: \qty{4}{\degreeCelsius} पर शुद्ध जल
का घनत्व \qty{1000.0}{kg/m^3} है, लवणता 35 वाले समुद्री जल का लगभग
\qty{1027.8}{kg/m^3} (आरेख पर पढ़ा गया)।
\end{solution}

\begin{solution}{exo:g10:layered-ocean:5}
प्लव \qty{1000}{m} पर बहता है, \qty{2000}{m} तक डूबता है, फिर रास्ते में
तापमान, लवणता (चालकता से) और दाब मापते हुए सतह तक उठता है;
सतह पर वह अपना प्रोफ़ाइल उपग्रह से भेजता है और फिर डूब जाता है।
प्लव धाराओं के साथ पूरे महासागर में बहते हैं और वर्षों तक, हर
ऋतु में, जहाज़ों के मार्गों से दूर काम करते हैं।
\end{solution}

\begin{solution}{exo:g10:layered-ocean:6}
$361.84\times10^6 \times 3.682 \approx \qty{1.33e9}{km^3}$, लगभग
\num{1.33} अरब घन किलोमीटर। $100/3682 \approx 0.027$: माध्य गहराई के
\qty{3}{\percent} से कम।
\end{solution}

\begin{solution}{exo:g10:layered-ocean:7}
$22.7 - 0.5 = \qty{22.2}{\degreeCelsius}$, जो \qty{15}{m} पर पहुँचता है: \omterm{def:g10:layered-ocean:mixed-layer}{मिश्रित
परत} लगभग \qty{15}{m} गहरी है। \qty{20}{m} और \qty{50}{m} के बीच:
$(21.8 - 17.3)/30 = \qty{0.15}{\degreeCelsius}$ प्रति मीटर।
\end{solution}

\begin{solution}{exo:g10:layered-ocean:8}
सूर्य केवल सतही परत को गरम करता है, और मौसमी गरम होना और ठंडा होना
केवल \omterm{def:g10:layered-ocean:mixed-layer}{मिश्रित परत} में फैलता है: \omterm{def:g10:layered-ocean:thermocline}{ताप-प्रवणता परत} के नीचे का जल
इनसे कटा रहता है। सर्दियों में सतह ठंडी होती है, अधिक घनी होकर नीचे के जल में
डूबती है, और तूफ़ान समुद्र को मथते हैं: मिश्रण गहराई तक जाता है। गर्मियों में
सूर्य एक हल्की सतही परत को गरम करता है जो ठंडे जल पर तैरती है; मंद
पवनें उसे नीचे नहीं धकेल पातीं, और \omterm{def:g10:layered-ocean:mixed-layer}{मिश्रित परत} पतली रहती है।
\end{solution}

\begin{solution}{exo:g10:layered-ocean:9}
\qty{10}{\degreeCelsius} के पास $0.78/0.18 \approx 4.3$ अंश;
\qty{0}{\degreeCelsius} के पास $0.78/0.06 = 13$ अंश। ठंडे ध्रुवीय समुद्रों में
घनत्व के लिए लवणता के छोटे परिवर्तन तापमान के परिवर्तनों से कहीं अधिक
मायने रखते हैं: वहाँ परतबंदी को लवणता नियंत्रित करती है।
\end{solution}

\begin{solution}{exo:g10:layered-ocean:10}
तापमान एक-सा है: गहरा जल इसलिए अधिक घना है कि वह
अधिक खारा है (33.3, जबकि 28.1)। उनके बीच की परत, जहाँ लवणता गहराई के साथ तेज़ी से
बदलती है, एक \omterm{def:g10:layered-ocean:thermocline}{लवण-प्रवणता परत} है; यहाँ वह \omterm{def:g10:layered-ocean:thermocline}{घनत्व-प्रवणता परत} भी है।
\end{solution}

\begin{solution}{exo:g10:layered-ocean:11}
नीचे धकेले जाने पर पार्सल अपना तापमान और लवणता, और इसलिए अपना घनत्व, बनाए रखता है।
घने जल से घिरा होने पर वह अपने द्वारा हटाए गए जल से हल्का होता है:
उत्प्लावन उसे वापस उसके अपने तल तक धकेल देता है। ऊपर धकेला गया पार्सल अपने
नए परिवेश से अधिक घना होगा और वापस डूब जाएगा। परतें स्थिर हैं।
\end{solution}

\begin{solution}{exo:g10:layered-ocean:12}
द्रव्यमान: $0.001 \times 1028.1 = \qty{1.0281}{kg}$। नया आयतन: $1.0281/1023.3
\approx \qty{0.0010047}{m^3}$, लगभग \qty{1004.7}{mL}: वह लगभग
\qty{4.7}{mL}, यानी \qty{0.47}{\percent}, फैला है। \qty{700}{m} की परत
$700 \times 0.0047 \approx \qty{3.3}{m}$ मोटी हो जाएगी।
\end{solution}

\begin{solution}{exo:g10:layered-ocean:13}
फ़रवरी में उत्तरी अटलांटिक में। वहाँ ऊपरी \qty{150}{m} का तापमान लगभग
एक-सा है, इसलिए घनत्व भी लगभग एक-सा: उसे मथने में थोड़ी ही ऊर्जा लगती है।
उष्णकटिबंधीय प्रशांत में \qty{50}{m} और \qty{200}{m} के बीच घनत्व $1026.0 - 1021.6 =
\qty{4.4}{kg/m^3}$ बढ़ता है: एक प्रबल
\omterm{def:g10:layered-ocean:thermocline}{घनत्व-प्रवणता परत}, और तूफ़ान को हल्के जल को कहीं अधिक घने जल में नीचे
धकेलना पड़ता।
\end{solution}

\begin{solution}{exo:g10:layered-ocean:14}
तटों से दूर उच्च पर्णहरित: प्रशांत में (और अटलांटिक में) भूमध्य रेखा के अनुदिश की पट्टी,
जहाँ व्यापारिक पवनें \omterm{def:g10:layered-ocean:upwelling}{उत्स्रवण} चलाती हैं; और उच्च
अक्षांश, जैसे उत्तरी अटलांटिक, जहाँ सर्दियों की \omterm{def:g10:layered-ocean:mixed-layer}{मिश्रित परत} गहराई तक
पहुँचती है और पोषक तत्त्वों को ऊपर लाती है। बहुत कम पर्णहरित: प्रबल रूप से \omterm{def:g10:layered-ocean:stratification}{स्तरित}
उपोष्ण वलयों के केंद्र, जहाँ कुछ भी गहरे जल को ऊपर नहीं लाता
और सूर्य-प्रकाशित परत के पोषक तत्त्व समाप्त हो चुके हैं।
\end{solution}

\begin{solution}{exo:g10:layered-ocean:15}
डूबना: $1000/(6 \times 3600) \approx \qty{0.046}{m/s}$, लगभग
\qty{4.6}{cm/s}। उठना: $2000/(10 \times 3600) \approx \qty{0.056}{m/s}$,
लगभग \qty{5.6}{cm/s}। $36\,000/2 = 18\,000$ मापन, जिनके बीच औसतन
$2000/18\,000 \approx \qty{0.11}{m}$ का अंतर है।
\end{solution}

\begin{solution}{pb:g10:layered-ocean:1}
\textbf{1.} वलय, $24.0 - 19.0 = \qty{5.0}{\degreeCelsius}$ से।

\textbf{2.} वलय (\qty{20.5}{\degree} दक्षिण, जबकि \qty{13.5}{\degree} दक्षिण)।
भूमध्य रेखा से अधिक दूर होने के कारण उसे अधिक ठंडा होना चाहिए: केवल अक्षांश
यह नहीं समझा सकते कि वह \qty{5}{\degreeCelsius} अधिक गरम क्यों है।

\textbf{3.} $24.0 - 0.5 = \qty{23.5}{\degreeCelsius}$, जो
\qty{50}{m} (23.6) और \qty{100}{m} (22.6) के बीच, लगभग $50 + 50 \times
0.1/1.0 = \qty{55}{m}$ पर पहुँचता है।

\textbf{4.} नहीं: तापमान सतह से ही लगातार गिरता है, पहले \qty{25}{m} में
\qty{1.6}{\degreeCelsius} और पहले \qty{100}{m} में \qty{5.3}{\degreeCelsius};
\omterm{def:g10:layered-ocean:thermocline}{ताप-प्रवणता परत} सतह से ही शुरू होती है।

\textbf{5.} \qty{200}{m} और \qty{300}{m} के बीच: $20.3 - 14.4 =
\qty{5.9}{\degreeCelsius}$ प्रति सौ मीटर।

\textbf{6.} पेरू के पास, \qty{50}{m} और \qty{100}{m} के बीच: अंतर्वेशन से लगभग
\qty{70}{m} (पूरा प्रोफ़ाइल \qty{65}{m} देता है)। वलय में,
\qty{200}{m} और \qty{300}{m} के बीच: लगभग $200 + 100 \times 5.3/5.9
\approx \qty{290}{m}$।

\textbf{7.} हाँ, दोनों स्तंभों के हर तल पर: हल्का जल घने
जल के ऊपर टिका है, और दोनों स्तंभ स्थिर हैं।

\textbf{8.} अधिक खारा जल अधिक घना होना चाहिए; फिर भी वलय का सतही जल
हल्का है: \qty{1024.71}{kg/m^3}, जबकि \qty{1025.08}{kg/m^3}।

\textbf{9.} वलय का जल लगभग 1.3 इकाई अधिक खारा है, जो लगभग
$1.3 \times 0.8 \approx \qty{1.0}{kg/m^3}$ जोड़ता है; वह \qty{5}{\degreeCelsius}
अधिक गरम है, जो इन तापमानों पर लगभग $5 \times 0.27 \approx \qty{1.4}{kg/m^3}$
घटाता है। तापमान जीतता है।

\textbf{10.} $1026.44 - 1025.08 = \qty{1.36}{kg/m^3}$।

\textbf{11.} $1025.30 - 1024.71 = \qty{0.59}{kg/m^3}$।

\textbf{12.} सतह से \qty{500}{m} तक घनत्व
$1026.92 - 1024.71 = \qty{2.21}{kg/m^3}$ बढ़ता है, जिसमें से केवल \qty{0.59}{kg/m^3}
\qty{200}{m} के ऊपर: इसका अधिकांश भाग, \qty{1.62}{kg/m^3}, \qty{200}{m} के नीचे है,
\qty{200}{m} और \qty{500}{m} के बीच की \omterm{def:g10:layered-ocean:thermocline}{ताप-प्रवणता परत} में।

\textbf{13.} पेरू के पास सतह से \qty{500}{m} तक घनत्व $1026.97 - 1025.08 =
\qty{1.89}{kg/m^3}$ बढ़ता है, जिसमें से
\qty{1.36}{kg/m^3} ऊपरी \qty{200}{m} में: इसका अधिकांश भाग सतह के पास है।

\textbf{14.} \qty{500}{m} पर दोनों स्थल लगभग एक जैसे हैं
(\qty{7.9}{\degreeCelsius} और \qty{7.1}{\degreeCelsius},
\qty{1026.97}{kg/m^3} और \qty{1026.92}{kg/m^3}): दोनों के नीचे वही ठंडा जल है;
अंतर यह है कि वह सतह के कितने पास आता है।

\textbf{15.} किनारे के अनुदिश भूमध्य रेखा की ओर बहती पवन
सतही जल को समुद्र की ओर धकेलती है; उसका स्थान लेने के लिए नीचे का जल तट के पास ऊपर
उठता है। \omterm{def:g10:layered-ocean:thermocline}{ताप-प्रवणता परत} का ठंडा जल ऊपर खिंच आता है, इसलिए \omterm{def:g10:layered-ocean:thermocline}{ताप-प्रवणता परत}
तट की ओर ऊपर झुकती है और लगभग सतह तक पहुँच जाती है।

\textbf{16.} पेरू के पास: \omterm{def:g10:layered-ocean:thermocline}{ताप-प्रवणता परत} के नीचे से आया \qty{13}{\degreeCelsius} से
\qty{15}{\degreeCelsius} का जल ऊपरी
\qty{100}{m} के भीतर है, पवन के मथने और प्रकाश दोनों की पहुँच में।
वलय में वह \qty{250}{m} से नीचे है, एक गरम, हल्की परत के नीचे जिसे पवन
उतनी गहराई तक नहीं मथ सकती।

\textbf{17.} \omterm{def:g10:layered-ocean:upwelling}{उत्स्रवण} रुक जाता: सतही जल धूप में गरम होता,
ऊपर एक गरम हल्की परत बन जाती, और सूर्य-प्रकाशित परत को पोषक तत्त्वों की
आपूर्ति घटती जाती।

\textbf{18.} पेरू का तट हल्के नीले से पीला है (बहुत पर्णहरित);
दक्षिण प्रशांत वलय का मध्य भाग मानचित्र के सबसे गहरे नीले क्षेत्रों में है
(बहुत कम)।

\textbf{19.} पादप-प्लवकों को प्रकाश के साथ-साथ पोषक तत्त्व भी चाहिए। वलय में
प्रबल \omterm{def:g10:layered-ocean:stratification}{स्तरीकरण} गहरे, पोषक-समृद्ध जल को सूर्य-प्रकाशित परत से बहुत नीचे
रखता है, और सतही परत के पोषक तत्त्व बहुत पहले समाप्त हो चुके हैं।

\textbf{20.} $1.36/0.59 \approx 2.3$: \textbf{ऊपरी \qty{200}{m} में घनत्व-वृद्धि
पेरू के पास वलय की तुलना में लगभग 2.3 गुना अधिक है}।
\omterm{def:g10:layered-ocean:thermocline}{घनत्व-प्रवणता परत}, अपने नीचे के ठंडे, पोषक-समृद्ध जल के साथ, पेरू के पास
ऊपरी \qty{200}{m} में उठा दी गई है, जबकि वलय में वह अधिक गहराई पर,
\qty{200}{m} के नीचे, है।
\end{solution}
